\begin{definition}[$N$-site Hilbert space]\label{def:nsite_space_parent}
    \lean{MPSTensor.NSiteSpace}
    \leanok
    For local dimension $d$ and chain length $N$, the $N$-site Hilbert space is
    \[
        \mathcal H_{N}^{(d)} := \{0,\ldots,d{-}1\}^{N} \to \C.
    \]
    This is the computational-basis realization used throughout the chapter for
    periodic-chain states.
\end{definition}

\section{\texorpdfstring{Local ground space $G_L(A)$}{Local ground space G\_L(A)}}\label{sec:ground_space_parent}

Fix an MPS tensor $A = \{A^i\}_{i=0}^{d-1}$ with $A^i \in \MN{D}$ and a block
length $L \ge 1$.
For a word $w = (i_1,\ldots,i_L)$, write
$A^w := A^{i_1}\cdots A^{i_L}$.
The boundary-condition parametrization is
\[
    \Gamma_L : \MN{D} \to (\{0,\ldots,d{-}1\}^L \to \C),
    \qquad
    \Gamma_L(X)(\sigma) := \tr\!(A^{\sigma} X),
\]
where $A^{\sigma}$ abbreviates
$A^{\sigma(0)}\cdots A^{\sigma(L-1)}$.

\begin{definition}[Ground-space map]\label{def:ground_space_map}
    \lean{MPSTensor.groundSpaceMap}
    \leanok
    \uses{def:eval_word}
    The map $\Gamma_L$ above is a $\C$-linear map
    \[
        \Gamma_L : \MN{D} \to (\{0,\ldots,d{-}1\}^L \to \C).
    \]
    In tensor-network notation,
    \begin{center}
        \TNGroundSpaceMap{A}{\sigma_1}{\sigma_L}{L}{X}
    \end{center}
    the upper chain denotes the word tensor $A^\sigma$ with open physical
    indices $\sigma_1,\ldots,\sigma_L$, and the lower red node denotes the boundary
    matrix $X$ inserted on the closing virtual bond before taking the trace.
    The chain length~$L$ is fixed by the surrounding formula, not by an extra
    label inside the diagram.
\end{definition}

\begin{lemma}[Evaluation formula]\label{lem:ground_space_map_apply}
    \lean{MPSTensor.groundSpaceMap_apply}
    \leanok
    \uses{def:ground_space_map}
    For every boundary matrix $X \in \MN{D}$ and every length-$L$ word $\sigma$,
    \[
        \Gamma_{L}(X)(\sigma) = \tr(A^{\sigma} X).
    \]
\end{lemma}

\begin{definition}[Local ground space]\label{def:ground_space_parent}
    \lean{MPSTensor.groundSpace}
    \leanok
    \uses{def:ground_space_map}
    The \emph{local ground space} is the image
    \[
        G_L(A) := \mathrm{im}(\Gamma_L)
        \subseteq (\{0,\ldots,d{-}1\}^L \to \C).
    \]
\end{definition}

\begin{lemma}[Dimension bound]\label{lem:ground_space_finrank_le}
    \lean{MPSTensor.groundSpace_finrank_le}
    \leanok
    \uses{def:ground_space_parent}
    For every $L$,
    \[
        \dim G_L(A) \le D^2.
    \]
\end{lemma}

\begin{proof}
    \leanok
    Since $G_L(A)$ is the range of a linear map from $\MN{D}$,
    \[
        \dim G_L(A) \le \dim \MN{D} = D^2.
    \]
\end{proof}

\begin{lemma}[Ambient-space dimension]\label{lem:nsite_space_finrank}
    \lean{MPSTensor.nSiteSpace_finrank}
    \leanok
    The above Hilbert-space realization satisfies
    \[
        \dim\!(\{0,\ldots,d{-}1\}^L \to \C)=d^L.
    \]
\end{lemma}

\begin{lemma}[Nontriviality criterion]\label{lem:ground_space_ne_top}
    \lean{MPSTensor.groundSpace_ne_top}
    \leanok
    \uses{lem:ground_space_finrank_le, lem:nsite_space_finrank}
    If $d^L > D^2$, then $G_L(A)$ is a proper subspace of the local Hilbert
    space:
    \[
        G_L(A) \neq (\{0,\ldots,d{-}1\}^L \to \C).
    \]
\end{lemma}

\begin{proof}
    \leanok
    If $G_L(A)$ were the whole space, its dimension would be $d^L$ by
    Lemma~\ref{lem:nsite_space_finrank};
    Lemma~\ref{lem:ground_space_finrank_le} gives the contradiction
    $d^L \le D^2$.
\end{proof}

\section{Parent interaction and chain Hamiltonian}\label{sec:parent_interaction_defs}

The canonical parent interaction at range $L$ is
\[
    h_L(A)=1-\Pi_{G_L(A)}=\Pi_{G_L(A)^\perp},
    \qquad
    \ker h_L(A)=G_L(A).
\]
The sources also allow any positive local interaction with this kernel; the
orthogonal projector is the canonical representative used here.

\begin{definition}[Hilbert-space copy of the local ground space]\label{def:ground_space_es}
    \lean{MPSTensor.groundSpaceES}
    \leanok
    \uses{def:nsite_space_parent, def:ground_space_parent}
    Let
    \[
        \iota_L :
        (\{0,\ldots,d{-}1\}^{L}\to\C)
        \simeq \ell^{2}(\{0,\ldots,d{-}1\}^{L})
    \]
    be the canonical computational-basis identification. The Hilbert-space copy
    of the local ground space is
    \[
        G_{L}^{\mathrm{ES}}(A):=\iota_L(G_L(A))
        \subseteq \ell^{2}(\{0,\ldots,d{-}1\}^{L}).
    \]
\end{definition}

\begin{lemma}[Membership in the Hilbert-space copy]\label{lem:mem_ground_space_es}
    \lean{MPSTensor.mem_groundSpaceES_iff}
    \leanok
    \uses{def:ground_space_es, def:ground_space_parent}
    For $v\in\ell^{2}(\{0,\ldots,d{-}1\}^{L})$,
    \[
        v\in G_L^{\mathrm{ES}}(A)
        \Longleftrightarrow
        \iota_L^{-1}v\in G_L(A).
    \]
\end{lemma}

\begin{proof}
    \leanok
    This is the defining equation
    $G_L^{\mathrm{ES}}(A)=\iota_L(G_L(A))$.
\end{proof}

\begin{definition}[Canonical parent interaction]\label{def:parent_interaction}
    \lean{MPSTensor.parentInteraction}
    \leanok
    \uses{def:ground_space_parent}
    The parent interaction at range $L$ is the canonical orthogonal projector
    \[
        h_L(A) := \Pi_{G_L(A)^\perp}=1-\Pi_{G_L(A)},
        \qquad
        \ker h_L(A)=G_L(A),
    \]
    on the local $L$-site space.
\end{definition}

\begin{lemma}[Parent interaction annihilates ground-space elements]%
    \label{lem:parent_interaction_apply_mem_ground_space}
    \lean{MPSTensor.parentInteraction_apply_mem_groundSpace}
    \leanok
    \uses{def:parent_interaction, def:ground_space_parent}
    For any $v \in G_L(A)$, we have $h_L(A)\,v = 0$.
\end{lemma}

\begin{proof}
    \leanok
    Since $h_L(A)$ is the orthogonal projector onto $G_L(A)^\perp$, it
    annihilates every element of $G_L(A)$.
\end{proof}

\begin{definition}[Restriction to a cyclic interval]\label{def:extract_window}
    \lean{MPSTensor.extractWindow}
    \leanok
    \uses{def:nsite_space_parent}
    For $\sigma \in \{0,\ldots,d{-}1\}^{N}$ and a starting site $i$, write
    \[
        \sigma|_{[i,i+L-1]}
        \in \{0,\ldots,d{-}1\}^{L}
    \]
    for the word defined by
    \[
        \bigl(\sigma|_{[i,i+L-1]}\bigr)(r)
        =\sigma(i+r \bmod N),
        \qquad 0\le r<L.
    \]
\end{definition}

\begin{definition}[Prescribing a cyclic interval]\label{def:replace_window}
    \lean{MPSTensor.replaceWindow}
    \leanok
    \uses{def:extract_window}
    If $L \le N$ and $\tau\in\{0,\ldots,d{-}1\}^{L}$, write
    \[
        \sigma^{[i,i+L-1]\leftarrow\tau}
    \]
    for the configuration defined by
    \[
        \sigma^{[i,i+L-1]\leftarrow\tau}(k)
        =
        \begin{cases}
            \tau(r),
                & k\equiv i+r \pmod N\text{ for some }0\le r<L,\\
            \sigma(k), & \text{otherwise}.
        \end{cases}
    \]
    The index $r$ in the first case is unique because $L\le N$.
\end{definition}

\begin{lemma}[Extracting a replaced window]\label{lem:extract_window_replace_window}
    \lean{MPSTensor.extractWindow_replaceWindow}
    \leanok
    \uses{def:extract_window, def:replace_window}
    If $L \le N$, then restricting to the same cyclic interval after prescribing
    its values recovers the inserted word:
    \[
        \left(\sigma^{[i,i+L-1]\leftarrow\tau}\right)|_{[i,i+L-1]}=\tau.
    \]
\end{lemma}

\begin{lemma}[Replacing an extracted window]\label{lem:replace_window_extract_window}
    \lean{MPSTensor.replaceWindow_extractWindow}
    \leanok
    \uses{def:extract_window, def:replace_window}
    If $L \le N$, then prescribing on a cyclic interval the values already
    carried by $\sigma$ leaves $\sigma$ unchanged:
    \[
        \sigma^{[i,i+L-1]\leftarrow\sigma|_{[i,i+L-1]}}=\sigma.
    \]
\end{lemma}
\begin{proof}\leanok
    Case-split on whether the offset $(k - i + N) \bmod N$ lies in $[0,L)$:
    inside the cyclic interval the prescribed value is the original value of
    $\sigma$; outside the interval both sides already agree with $\sigma$.
\end{proof}

For a periodic chain of length $N$, the parent Hamiltonian is the sum of
translated copies of the local interaction.

\begin{definition}[Translated local term]\label{def:local_term_parent}
    \lean{MPSTensor.localTerm}
    \leanok
    \uses{def:parent_interaction, def:replace_window, def:extract_window}
    For each site index $i \in \{0,\ldots,N{-}1\}$, the translated local term is
    determined by
    \[
        (h_i\psi)(\sigma)
        =
        \begin{cases}
        \left[
            h_L(A)
            \bigl(\tau\mapsto
                \psi(\sigma^{[i,i+L-1]\leftarrow\tau})\bigr)
        \right]\!\bigl(\sigma|_{[i,i+L-1]}\bigr), & L \le N,\\
        0, & L > N.
        \end{cases}
    \]
\end{definition}

\begin{definition}[Parent Hamiltonian]\label{def:parent_hamiltonian}
    \lean{MPSTensor.parentHamiltonian}
    \leanok
    \uses{def:local_term_parent}
    The chain Hamiltonian is
    \[
        H_N(A,L) := \sum_{i=0}^{N-1} h_i.
    \]
\end{definition}

\begin{definition}[Frustration-free condition]\label{def:is_frustration_free}
    \lean{MPSTensor.IsFrustrationFree}
    \leanok
    \uses{def:local_term_parent}
    A vector $\psi$ is a frustration-free ground state of the parent Hamiltonian
    when each local interaction annihilates it:
    \[
        \forall\, i \in \{0,\ldots,N{-}1\},\quad h_i \psi = 0.
    \]
\end{definition}

\begin{lemma}[MPV window membership]\label{lem:mpv_window_mem_ground_space}
    \lean{MPSTensor.mpv_window_mem_groundSpace}
    \leanok
    \uses{def:ground_space_parent, def:mpv}
    For any window of $L$ consecutive sites starting at position~$i$ in an
    $N$-site periodic chain ($L \le N$), the restriction of the MPV state to
    that window lies in $G_L(A)$.
\end{lemma}

\begin{proof}
    \leanok
    The boundary matrix is the product of $A$-matrices on the complement sites.
    Trace cyclicity rotates the full $N$-site product so that window
    indices come first, matching the ground-space map definition.
\end{proof}

\begin{lemma}[Local term annihilates the MPV]\label{lem:local_term_annihilates_mpv}
    \lean{MPSTensor.localTerm_annihilates_mpv}
    \leanok
    \uses{def:local_term_parent, def:mpv, lem:mpv_window_mem_ground_space,
        lem:parent_interaction_apply_mem_ground_space}
    For each site $i$, $h_i\,V^{(N)}(A) = 0$.
\end{lemma}

\begin{proof}
    \leanok
    Combines MPV window membership with the fact that the parent
    interaction annihilates ground-space elements.
\end{proof}

\begin{lemma}[Parent Hamiltonian annihilates the MPV]\label{lem:parent_hamiltonian_annihilates}
    \lean{MPSTensor.parentHamiltonian_annihilates}
    \leanok
    \uses{def:parent_hamiltonian, def:mpv, lem:local_term_annihilates_mpv}
    For every $N$, the MPV state satisfies
    \[
        H_N(A,L)\,V^{(N)}(A) = 0.
    \]
\end{lemma}

\begin{proof}
    \leanok
    Sum of zeros, since each local term annihilates the MPV.
\end{proof}

\begin{lemma}[Frustration-freeness of the MPV]\label{lem:parent_hamiltonian_ff}
    \lean{MPSTensor.parentHamiltonian_frustrationFree}
    \leanok
    \uses{def:is_frustration_free, def:mpv, lem:local_term_annihilates_mpv}
    The MPS vector is a frustration-free ground state: for every site $i$,
    \[
        h_i V^{(N)}(A)=0.
    \]
\end{lemma}

\begin{proof}
    \leanok
    Immediate from Lemma~\ref{lem:local_term_annihilates_mpv}.
\end{proof}
